\section*{Task 2: Main stratification theorem}

\paragraph{Goal.}
State the main result of the paper in the final weight-space language.

\paragraph{Main theorem.}
For the fixed singular subset \(S\subseteq \{1,\dots,r_{\max}\}\) with
\(|S|=r\), the rank-\(r\) critical locus admits the disjoint decomposition
\[
\mathcal C_{d,S}^r
=
\bigsqcup_{\rho\in \widetilde{\mathcal R}_{d,S}^{\,r}}
\mathcal C_{d,S}^r[\rho],
\qquad
\mathcal C_{d,S}^r[\rho]
:=
\mathcal C_{d,S}^r\cap \Gamma_{d,S}^r[\rho].
\]
Each nonempty piece \(\mathcal C_{d,S}^r[\rho]\) is a single \(H_S\)-orbit.
Equivalently,
\[
H_S\backslash \mathcal C_{d,S}^r
\xrightarrow{\sim}
\widetilde{\mathcal R}_{d,S}^{\,r}.
\]

\paragraph{Interpretation.}
This theorem gives the discrete combinatorial label of a saddle stratum:
every critical point of rank \(r\) has a unique admissible cyclic rank pattern,
and two critical points lie in the same \(H_S\)-orbit if and only if they have
the same lifted cyclic rank pattern.

\paragraph{Proof.}
The cyclic locus decomposition from Conjecture 3 gives
\[
\Gamma_{d,S}^r
=
\bigsqcup_{\rho\in \widetilde{\mathcal R}_{d,S}^{\,r}}
\Gamma_{d,S}^r[\rho],
\]
with each nonempty piece a single \(G_d^{W_0}\)-orbit. The first-conjecture
orbit correspondence identifies each such ambient orbit with exactly one
\(H_S\)-orbit in the critical locus. Intersecting the ambient decomposition with
\(\mathcal C_{d,S}^r\) yields the stated partition.

\paragraph{Why this is the main theorem.}
The theorem turns the degenerate critical geometry of deep linear networks into
a finite combinatorial classification problem. The role of the cyclic rank
pattern is exactly to provide the discrete identity of a saddle stratum inside a
non-isolated critical locus.

\paragraph{Status.}
Completed. This is the main critical-locus parametrization theorem needed for
the paper.
